\documentclass[11pt,oneside]{article}
\input{../../../Header/packages.sty}
\input{../../../Header/mathexpressions.sty}
\usepackage{fancyhdr}
\pagestyle{fancy}
\fancyhf{}
\lhead{Lektionsplanung: Der Kathetensatz}
\rhead{Jonas Guler (Gu)}
\rfoot{\thepage}
\renewcommand{\headrulewidth}{0.4pt}
\setlength{\headheight}{14pt}
\setlength{\parindent}{0pt}
\setlength{\parskip}{4pt}
\linespread{1.15}

\begin{document}
{\Large\textbf{Lektionsplanung: Der Kathetensatz}}\\[2mm]
\begin{tabularx}{\textwidth}{@{}lX}
	Thema: & Sätze im rechtwinkligen Dreieck, Kathetensatz (Satz des Euklid)\\
	Dauer: & 1 Lektion à 45 Minuten\\
	Einbettung: & nach dem Satz des Pythagoras (Skript \textit{Sätze im rechtwinkligen Dreieck}, Abschnitt 3), vor dem Höhensatz\\
	Material: & Arbeitsblatt \texttt{kathetensatz.pdf}, Taschenrechner, Geodreieck und Zirkel\\
\end{tabularx}

\section*{Lernziele}
Die Lernenden \ldots
\begin{itemize}[itemsep=0pt]
	\item entdecken an Zahlenbeispielen den Zusammenhang $a^2=c\cdot q$ bzw.\ $b^2=c\cdot p$ und formulieren ihn als Vermutung.
	\item können den Kathetensatz in Worten und als Formel angeben und am Flächenbild (Kathetenquadrat = Rechteck) erklären.
	\item verstehen eine Herleitung des Kathetensatzes, die nur den Satz des Pythagoras benutzt, und wissen, dass umgekehrt Pythagoras aus dem Kathetensatz folgt.
	\item berechnen mit dem Kathetensatz fehlende Grössen im rechtwinkligen Dreieck, insbesondere die Hypotenusenabschnitte.
\end{itemize}

\textbf{Vorwissen:} Satz des Pythagoras, Bezeichnungen $a,b,c,h_c,p,q$ (Skript Abschnitt 1), Berechnung von $h_c$ über den Flächeninhalt (Skript Aufgabe 2.1 Nr.\,6).

\textbf{Achtung Notation:} Im Skript liegt $p$ bei $A$ (also an der Kathete $b$) und $q$ bei $B$. Es gilt deshalb $a^2=c\cdot q$ und $b^2=c\cdot p$. In \textit{Geometrie 2} (Niederberger) ist es umgekehrt ($a^2=c\cdot p$). Das Arbeitsblatt folgt konsequent dem Skript.

\section*{Ablauf}
\renewcommand{\arraystretch}{1.3}
\begin{tabularx}{\textwidth}{|>{\raggedright\arraybackslash}p{11mm}|>{\raggedright\arraybackslash}p{27mm}|X|>{\raggedright\arraybackslash}p{22mm}|}\hline
	\textbf{Zeit} & \textbf{Phase} & \textbf{Inhalt und Aktivität} & \textbf{Sozialform}\\\hline
	0'--2' & Einstieg & Begrüssung, Lektionsziel nennen: \textit{Heute finden wir mit Pythagoras einen neuen Satz im rechtwinkligen Dreieck.} Arbeitsblatt austeilen, Bezeichnungen an der Figur auf S.\,1 kurz zeigen ($p$ an $b$, $q$ an $a$). & Plenum\\\hline
	2'--14' & Schüleraktivität (Entdeckung) & \textbf{Aufgabe 1.1, Teil A:} $c$, $h_c$, $p$, $q$ für $a=15$, $b=20$ berechnen, Tabelle ausfüllen und Vermutung formulieren. \textbf{Teil B} (Beweis) für schnelle Paare. Lehrperson geht herum; häufige Hürde ist $h_c$: auf den Flächeninhalt $\frac{a\cdot b}{2}=\frac{c\cdot h_c}{2}$ hinweisen. & Partnerarbeit\\\hline
	14'--26' & Theorieinput & Vermutungen sammeln (2--3 Paare). \textbf{Satz 2.1} an der Wandtafel mit Flächenbild entwickeln: farbiges Kathetenquadrat und gleichfarbiges Rechteck. \textbf{Beweis} gemeinsam mit den drei Pythagoras-Gleichungen (evtl.\ von einem Paar aus Teil B vorstellen lassen). \textbf{Bemerkungen:} Pythagoras folgt aus dem Kathetensatz (Euklid), Formeln $q=\frac{a^2}{c}$, $p=\frac{b^2}{c}$, Merksatz in Worten. \textbf{Beispiel 1} gemeinsam lösen (ca.\ 3'). & Plenum, Lehrervortrag mit Fragen\\\hline
	26'--44' & Übungen & \textbf{Pflicht:} Aufgabe 3.1 a)--d), 3.2, 3.3.\newline \textbf{Zusatz:} 3.1 e)--f), 3.4 (Konstruktion mit Thaleskreis), 3.5 $\bigstar$ (Raumdiagonale). Lösungen liegen bei der Lehrperson auf bzw.\ werden ab Min.\ 38 aufgelegt. & Einzel- oder Partnerarbeit\\\hline
	44'--45' & Abschluss & Merksatz von einer Schülerin oder einem Schüler wiederholen lassen: \textit{Kathete$^2$ = Hypotenuse $\cdot$ anliegender Abschnitt.} Hausaufgabe: Pflichtaufgaben fertig lösen. & Plenum\\\hline
\end{tabularx}

\section*{Tafelbild (Vorschlag)}
\begin{minipage}{0.48\linewidth}
	\textbf{Kathetensatz (Satz des Euklid)}\\
	Kathetenquadrat = Rechteck aus Hypotenuse und anliegendem Hypotenusenabschnitt:
	\[a^2=c\cdot q,\qquad b^2=c\cdot p\]
	\textbf{Beweis:} (I) $a^2=h^2+q^2$, (II) $b^2=h^2+p^2$, (III) $a^2+b^2=p^2+2pq+q^2$.\\
	(I)$-$(II): $a^2-b^2=q^2-p^2$, dazu (III):\\ $2a^2=2q^2+2pq=2q(p+q)=2qc$.
\end{minipage}
\hfill
\begin{minipage}{0.45\linewidth}
	\centering
	\begin{tikzpicture}[x=1cm, y=1cm, font=\footnotesize, scale=0.42]
		\filldraw[fill=myblue, fill opacity=0.25, draw=black] (0,0) -- (3.2,2.4) -- (0.8,5.6) -- (-2.4,3.2) -- cycle;
		\filldraw[fill=orange, fill opacity=0.25, draw=black] (3.2,2.4) -- (5,0) -- (7.4,1.8) -- (5.6,4.2) -- cycle;
		\filldraw[fill=myblue, fill opacity=0.25, draw=black] (0,0) -- (3.2,0) -- (3.2,-5) -- (0,-5) -- cycle;
		\filldraw[fill=orange, fill opacity=0.25, draw=black] (3.2,0) -- (5,0) -- (5,-5) -- (3.2,-5) -- cycle;
		\draw[line width=1pt] (0,0) -- (5,0) -- (3.2,2.4) -- cycle;
		\draw[dashed] (3.2,2.4) -- (3.2,0);
		\draw (-0.6,2.6) node {$b^2$};
		\draw (5.3,2.1) node {$a^2$};
		\draw (1.6,-2.5) node {$c\, p$};
		\draw (4.1,-2.5) node {$c\, q$};
	\end{tikzpicture}
\end{minipage}

\section*{Mögliche Schwierigkeiten und Differenzierung}
\begin{itemize}[itemsep=0pt]
	\item \textbf{Zuordnung $p$/$q$:} Viele verwechseln, welcher Abschnitt zu welcher Kathete gehört. Immer mit Skizze arbeiten und den Merksatz in Worten (\textit{anliegender} Abschnitt) betonen.
	\item \textbf{Rechter Winkel nicht bei $C$:} In späteren Aufgaben bewusst thematisieren (analog zu Pythagoras im Skript).
	\item \textbf{Schwächere Lernende:} In Aufgabe 1.1 nur das zweite Dreieck (Werte gegeben) ausfüllen lassen und direkt die Vermutung formulieren.
	\item \textbf{Stärkere Lernende:} Teil B von Aufgabe 1.1, Aufgaben 3.4 und 3.5; zusätzlich: Kehrsatz formulieren (\textit{Gilt $a^2=c\cdot q$, ist das Dreieck dann rechtwinklig?}).
	\item \textbf{Ausblick nächste Lektion:} Höhensatz $h^2=p\cdot q$ (Geometrie 1 Aufgaben, Kap.\,5.2 Nr.\,40).
\end{itemize}

\newpage
\section*{Lösungen}
\textbf{Aufgabe 1.1}
\begin{itemize}[itemsep=0pt]
	\item[a)--c)] $c=\sqrt{15^2+20^2}=25$, $h_c=\frac{15\cdot 20}{25}=12$, $q=\sqrt{15^2-12^2}=9$, $p=\sqrt{20^2-12^2}=16$.
	\item[d)] Dreieck 1: $a^2=225=c\cdot q=25\cdot 9$, $b^2=400=c\cdot p=25\cdot 16$.\\ Dreieck 2: $a^2=36=10\cdot 3.6$, $b^2=64=10\cdot 6.4$.
	\item[e)] Vermutung: $a^2=c\cdot q$ und $b^2=c\cdot p$.
	\item[f)--g)] (I) $a^2=h_c^2+q^2$, (II) $b^2=h_c^2+p^2$, (III) $a^2+b^2=(p+q)^2=p^2+2pq+q^2$.\\
	(I)$-$(II): $a^2-b^2=q^2-p^2$. Addiert zu (III): $2a^2=2q^2+2pq=2q(p+q)=2qc$, also $a^2=c\cdot q$. Analog (oder mit (III)) $b^2=c\cdot p$.
\end{itemize}

\textbf{Beispiel 1:} $q=\frac{a^2}{c}=\frac{36}{9}=4$\,cm, $p=c-q=5$\,cm, $b^2=c\cdot p=45$, $b=\sqrt{45}=3\sqrt5\approx 6.71$\,cm. Kontrolle: $36+45=81=9^2$.

\textbf{Aufgabe 3.1}
\begin{enumerate}[label=\alph*), itemsep=0pt]
	\item $c=5$\,cm, $q=\frac{16}{5}=3.2$\,cm, $p=\frac{9}{5}=1.8$\,cm.
	\item $p=\frac{144}{16}=9$\,cm, $q=16-9=7$\,cm.
	\item $c=10$\,cm, $a=\sqrt{10\cdot 8}=4\sqrt5\approx 8.94$\,cm, $b=\sqrt{10\cdot 2}=2\sqrt5\approx 4.47$\,cm.
	\item $c=\frac{b^2}{p}=\frac{100}{8}=12.5$\,cm, $q=4.5$\,cm, $a=\sqrt{12.5\cdot 4.5}=7.5$\,cm.
	\item $q=\frac{25}{13}\approx 1.92$\,cm, $p=\frac{144}{13}\approx 11.08$\,cm, $b=12$\,cm.
	\item $c=\frac{12}{2}=6$\,cm, $p=4$\,cm, $b=\sqrt{24}=2\sqrt6\approx 4.90$\,cm.
\end{enumerate}

\textbf{Aufgabe 3.2:} Das Dreieck $\triangle ABC$ ist rechtwinklig bei $B$ mit Hypotenuse $\overline{AC}=25$\,cm. Kathetensatz: $\overline{AE}=\frac{20^2}{25}=16$\,cm, $\overline{CE}=\frac{15^2}{25}=9$\,cm, $\overline{BE}=\sqrt{20^2-16^2}=12$\,cm.

\textbf{Aufgabe 3.3:} $a^2=c\cdot q$ mit $a=2q$: $4q^2=10q$, also $q=2.5$\,cm, $a=5$\,cm, $p=7.5$\,cm, $b=\sqrt{10\cdot 7.5}=5\sqrt3\approx 8.66$\,cm.

\textbf{Aufgabe 3.4:} Das Dreieck $\triangle ABC$ ist nach Thales rechtwinklig bei $C$, $H$ ist der Höhenfusspunkt. Kathetensatz: $\overline{BC}^2=c\cdot q=4.5\cdot 2.5=11.25$\,cm$^2$, also $\overline{BC}\approx 3.35$\,cm.

\textbf{Aufgabe 3.5:} Dreieck $\triangle BDH$ ist rechtwinklig bei $D$ mit $\overline{BD}=10$\,cm, $\overline{DH}=7.5$\,cm, $\overline{BH}=12.5$\,cm. $\overline{BP}=\frac{10^2}{12.5}=8$\,cm, $\overline{HP}=\frac{7.5^2}{12.5}=4.5$\,cm, also $\overline{BP}:\overline{HP}=16:9$.

\section*{Quellen und Inspiration}
\begin{itemize}[itemsep=0pt]
	\item Geometrie 2 (Niederberger), Kap.\ \textit{Der Kathetensatz}, S.\,80--82 (Satz, Flächenverwandlung, Aufgaben 1, 4, 5, 7).
	\item Geometrie 1 Aufgaben, Kap.\,5.2 \textit{Höhensatz und Kathetensatz}, Nr.\,43--47.
\end{itemize}
\end{document}
